\begin{Def}
A Projective Plane is a Linear Space 
\((P,L)\) such that 
\begin{Axiom}(P1)
    Every Line has at least 3 distinct points.
\end{Axiom}
and 
\begin{Axiom}(P2)
    Any two lines meet in at least one point
\end{Axiom}
\end{Def}
\begin{Lemma}
Let \((P,L)\) be a projective space. Then for every point there are at least 3 lines through 
every point.
\end{Lemma}
\begin{proof}
Let \(p\in P\) be a point.
Since there are the non-collinear points, there are at least two distinct lines 
through \(p\). Let \(\ell,m\) be lines through \(p\) and pick \(A\in \ell\) and \(B\in m\) different 
from \(p\). Let \(o\) be the line through \(A\) and \(B\). Since 
\(k_o\geq 3\) there must be a point \(Q\in o\) distinc from \(A\) and \(B\)
Then \(Q\not \in \ell\) because otherwise \(o\) and \(\ell\) contain the distinct 
points \(A,Q\) so \(o=\ell\) which implies \(B\in \ell\).
Similarly \(Q\not in m\).
Then the line through \(Q\) and \(p\) is a third line through \(p\).


\end{proof}
\begin{Def}
Let \((P,L)\) be a projective space. Let \(P^\vee=L\) and 
for \(p\in P\) define \(p^\vee=\{\ell\in L\mid p\in \ell\}\) and let 
\(L^\vee=\{p^\vee\mid p\in P\}\).
Then \((P^\vee,L^\vee)\) is a projective space, called the dual space to \(P,V\) and 
\(k_{p^\vee}=r_p\) and \(k_\ell=r_\ell\).
\end{Def}
\begin{proof}
Let \(\ell,m\in P^\vee\) with \(\ell\neq m\). Then let \(p\) be the intersection point 
of \(m,\ell\). Then \(\ell,m\in p^\vee\). this is unique because if \(\ell,m\in q^\vee\) then
\(q\in \ell\cap m=\{p\}\).
Let \(p^\vee\in L^\vee\). Since \(r_p\geq 2\) we have that there are \(\ell,m\in p^\vee\) with \(\ell\neq m\).
Let \(A,B,C\in P\) be three non-collinear points.
Then \(\overline{AB},\overline{AC},\overline{BC}\in P^\vee\) are thre non-collinear points because 
if they were all in \
then all thre lines would have two points in common with each other, making them all equal contradicting
the fact that they are non-collinear.
now let \(p^\vee\in L^\vee\) for some \(p\in P\). Since 
\(r_p\geq 2\) we have that \(|p^\vee|\geq 3\).
Let \(p^\vee,q^\vee\in L^\vee\).
Since there is a line \(\ell\) through \(p\) and \(q\) we have \(\ell\in p^\vee\cap q^\vee\).
Thus \((P^\vee,L^\vee)\) is projective.

\end{proof}
\begin{Satz}
Let \((P,L)\) be a finite Linear Space such that \(k_\ell\geq 3\) for all \(\ell\in L\).
Then it is equivalent:
\begin{enumerate}
    \item \((P,L)\) is projective.
    \item \(|P|=|L|\).
\end{enumerate}
In this case \(k_\ell=r_{p}=q+1\) for all \(\ell\in L\) and \(p\in P\) and 
\(|P|=q^2+q+1\).
\end{Satz}
\begin{proof}
If \((P,L)\) is projective, then \(P^\vee,L^\vee\) is projective and 
\[|P|\leq |L|=|P^\vee|\leq |L^vee|\] 
Since there is a bijektion \(P\to L^\vee, p\mapsto p^\vee\) we have 
\[|P|\leq |L|\leq |P|\] such that 
\(|P|=|L|\).
If \(|P|=|L|\) then Satz applies to show that \((P,L)\) is projective and the formulas follow.
\end{proof}
\begin{Def}[Affine Plane]
    An affine Plane is a Linear Space \((P,L)\) satisfiyng \axiomref{Axiom:P'}, .ie
    For every line \(\ell\in L\) and every point \(p\in P\) with \(p\not\in\ell\) there is exactly 
    one line \(m\in\ell\) with \(p\in m\) and \(m\parallel \ell\).

\end{Def}
\begin{Satz}
Let \((P,L)\) be an Affine Plane.
Let \(\ell,m\in L\) be two lines. Then there is a bijection 
\(f\colon \ell\to m\).
\end{Satz}
\begin{proof}
Assume \(\ell\neq m\).
choose \(A\in \ell\setminus n\) and \(B\in m\setminus\ell\).
Define the function 
\(f\colon \ell\to m\) sending \(A\to B\) and a point \(X\neq A\) to the intersection of 
\(m\) with the unique line \(o_X\) parallel to \(\overline{AB}\) containing \(X\).
This must meet \(m\) because otherwise \(\ell=o_X\) by the uniqeness of parallel lines through \(X\)
but then \(A\in o_X\) ao \(o_X\) cannot be parallel to \(\overline{AB}\).
Thus \(f\) is well-defined.
SImilarly we define \(g\colon m\to \ell\) and it is easy to check that they are 
mutually inverse.
\end{proof}
\begin{Satz}
Let \((P,L)\) be a finite Affine Plane such that there is a line with \(n\) points.
Then there are exactly \(n+1\) distinct lines through every point and 
\(|P|=n^2\).
\end{Satz}
\begin{proof}
Let \(p\in P\) and \(\ell\in L\) with \(p\not \in L\).
Then for \(A,B\in\ell\) with \(A\neq B\) the line \(m\) through \(A\) and \(p\) and the line 
\(o\) through
\(B\) and \(p\) are distinct, because otherwise \(A,B\in o\) so \(o=\ell\) and \(p\in \ell\).
Thus there are at least \(n\) different lines through \(p\). There is one line through \(p\) parallel
to \(\ell\). so there are at least \(n+1\) lines through \(p\).
There cannot be more, because if we take any line through \(p\) it is either parallel to \(\ell\) making 
it the unique line through \(p\) or it intersects with \(\ell\) in which case it is 
\(\overline{Ap}\) from above.
so \(r_p=n+1\).
By ??? we have 
\[n+1=\dfrac{|P|-1}{n-1}\] so that 
\[|P|=n^2\].
\end{proof}
\begin{Satz}
Let \((P,L)\) be a projective space.
Let \(\ell\in L\) be a line. For another line \(m\in L\setminus\{\ell\}\) define 
\(p_m\) as the intersection of \(m\) and \(\ell\).
Let \(P'=P\setminus\{ p\mid p\in \ell\}\) and 
\(L'=\{m\setminus\{p_m\}\mid \ell\neq m\in L\}\). 
Then \((P',L')\) is an affine space. Every affine space comes from a projective space in that way.
\end{Satz}
\begin{proof}
    we check only \axiomref{Axiom:P'}.
    Let \(m\in L\) and \(p\) a point not on \(\ell\) and not on \(m\).
    let \(o\in L\) be the line through \(p_m\) and \(p\).
    Then \(m\setminus\{p_m\}\) and \(o\setminus\{p_m\}\) are parallel with \(p\in o\).
    On the other hand, if \((P',L')\) is an affine space. 
    Let for \(\ell\in L\) \[p_\ell=[\ell]\] the equivalence class of \(\ell\) from parallelism.
    Then define \(P=P\cup\{[\ell]\mid \ell\in L\}\) and \(\ell_\infty=\{[\ell]\mid \ell\in L'\}\)
    \[L=\{\ell\cup\{[\ell]\}\mid \ell\in L'\}\cup\{\ell_\infty\}\] 
    Then \((P,L)\) is a projective space and removing \(\ell_\infty\) as above recovers \((P',L')\).
    We only show that two lines in \((P,L)\) meet.
    If \(m,\ell\in L'\) then they either meet in a point in \(P'\) or they are parallel 
    so that
    \(m\cup\{[m]\}\) and \(\ell\cup\{[\ell]\}\) meet in \([m]=[\ell]\).
    \(\ell\cup\{[\ell]\}\) and \(\ell_\infty\) meet in \([\ell\]\).
\end{proof}